\section{Introduction}

Conventional experimentation in materials discovery and chemical sciences remains slow and resource-intensive, relying on sequential cycles of human intuition, hypothesis-driven testing, and trial-and-error optimization. As a result, the pace of discovery often falls short of the growing demand for advanced materials and sustainable chemical processes \cite{lee2026toward, bayley2024autonomous}. In response, self-driving labs (SDLs) and materials acceleration platforms (MAPs) have emerged as a new transformative paradigm of accelerated scientific discovery, integrating automation, artificial intelligence (AI), and algorithmic decision-making into closed-loop workflows to perform autonomous experiments with minimal human intervention \cite{lee2026toward, bayley2024autonomous, hickman2023self, hase2019next}. Within these systems, AI models are continuously updated using real-time experimental feedback to capture the behavior of the underlying system and guide the selection of subsequent experiments \cite{flores2020materials,roch2020chemos, sim2024chemos}. This closed-loop cycle, composed of experimentation, measurement, model refinement, and decision-making, enables data-efficient progress toward predefined objectives such as discovering new materials \cite{szymanski2023autonomous}, optimizing functional properties \cite{macleod2020self}, or improving reaction performance \cite{burger2020mobile}. In addition, the effectiveness of SDLs depends not only on individual components but also critically on the design, coordination, and adaptability of the overall experimental workflow \cite{hysmith2024future}.

Active learning has consequently become a central algorithmic component of SDLs, as closed-loop experimentation fundamentally requires selecting informative experiments under uncertainty. In a typical active-learning workflow, a surrogate model is iteratively updated, and its predictions and uncertainties are combined through an acquisition function to recommend experiments that balance exploration and exploitation \cite{lookman2019active}. In materials and chemical sciences, Bayesian optimization (BO), particularly multi-objective Bayesian optimization (MOBO), has emerged as a well-established realization of this framework, owing to the suitability of probabilistic surrogate models and exploration-exploitation strategies for settings characterized by expensive, noisy, and data-scarce experiments. 

Recent studies have demonstrated and benchmarked BO and MOBO across diverse discovery tasks, including multi-property materials design, establishing them as foundational tools for data-efficient experimentation \cite{liang2021benchmarking,wang2022benchmarking,arroyave2022perspective,adesiji2026benchmarking, desimpel2026bayesian, padhy2024experimentally}. In practice, these approaches are typically deployed in sequential or batched settings, where multiple experiments are proposed at each iteration to improve throughput, further increasing the complexity of decision-making.

Despite these advances, most implementations of BO in autonomous experimentation rely on fixed or manually tuned acquisition strategies \cite{de2021greed}. Such approaches implicitly assume stable experimental conditions and consistent objectives, whereas real-world experimental campaigns are often subject to changing constraints, including evolving design goals, limited budgets, and time-critical decisions. Under these conditions, the optimal balance between exploration and exploitation is inherently non-stationary, and fixed strategies can lead to inefficient allocation of experimental effort \cite{hoffman2011portfolio}. While BO provides statistically grounded mechanisms for selecting candidate experiments, it does not explicitly address the higher-level problem of adaptive policy selection, i.e., determining how the exploration–exploitation trade-off should evolve throughout the course of a campaign. This limitation can be viewed as a separation between local decision-making, governed by acquisition functions, and global decision-making, which governs how experimental effort is allocated over time \cite{armesto2023acquisition}. The absence of mechanisms to bridge this separation becomes particularly pronounced in batched and multi-objective settings, where trade-offs must be managed across competing objectives, parallel experimental allocations, and finite resource budgets\cite{Daulton2020}.

In parallel with these developments, large language models (LLMs) have rapidly gained traction in chemical and materials research, primarily as interfaces for arduous tasks such as literature mining and knowledge extraction, planning, and tool-mediated assistance \cite{lei2024materials,jablonka202314,ramos2025review, jiang2025applications,dagdelen2024structured,katzer2025towards}. However, cautious assessments suggest that generic LLMs are not yet reliable domain reasoners and that effective deployment requires stronger grounding in domain-specific data, multimodal inputs, and structured interaction with expert tools \cite{miret2025enabling}. This perspective is supported by systems such as ChemCrow \cite{m2024augmenting}, which show that LLMs become substantially more useful when orchestrating search, safety, molecular, and reaction tools. CLAIRify \cite{yoshikawa2023large} and Coscientist \cite{boiko2023autonomous} demonstrate that language models can already participate in planning and execution layers of closed research workflows. Collectively, these studies suggest that LLMs are most effective as domain-aware coordinators that integrate prior knowledge, computational tools, laboratory systems, and human intent. This, in turn, raises the question of whether such systems can be extended to play structured roles in experimental decision-making.

In this context, a more concrete question is whether an LLM agent can dynamically manage a BO experiment, much like a human experimenter. Motivated by this perspective, recent work has begun to explore the role of LLMs in BO. For example, in chemistry, BoChemian \cite{rankovic2023bochemian} uses LLM embeddings of textual reaction procedures as feature representations for BO, whereas Kristiadi \etal \cite{kristiadi2024sober} takes a more cautious view in the discovery of molecular materials, showing that LLM-derived representations aid principled BO only when the underlying model is specific to the chemistry or adapted with domain data. In more general BO settings, LLAMBO \cite{liu2024large} leverages natural-language descriptions of prior evaluations for zero-shot warm starting, surrogate enhancement, and candidate sampling, highlighting the potential of language-based priors to influence early-stage search. 

More recent work has shifted from end-to-end proposal generation toward higher-level control. Gupta \emph{et al.} \cite{gupta2025llms} show that the open-source instruction-tuned LLM agents they evaluate are largely insensitive to experimental feedback and are consistently outperformed by linear bandits and Gaussian-process-based methods, motivating hybrid frameworks that separate prior-based reasoning from posterior-updating acquisition policies. In contrast, LMABO \cite{ngo2026adaptive} uses an LLM as a zero-shot strategist that selects among acquisition functions based on structured summaries of the optimization state, suggesting a more constrained and potentially more effective role for LLMs in guiding optimization. 

The findings described above point to a more structured role for LLMs—as bounded meta-decision layers that guide, rather than replace, established optimization frameworks. However, evidence for such roles in multi-objective optimization settings remains limited, particularly in scenarios where experimental constraints evolve dynamically and require adaptive reallocation of resources across competing objectives. This gap highlights the need for approaches that explicitly address decision-making at the level of policy selection, rather than candidate generation alone.

In this work, we address this gap by investigating whether an LLM-based agent can act as a high-level controller for batched MOBO. Specifically, we consider the problem of dynamically allocating experimental effort between exploration and exploitation under changing resource constraints, where the optimal strategy is inherently non-stationary. Our framework retains conventional MOBO components for surrogate modeling and batch acquisition, while introducing a separate decision layer that selects among structured allocation policies based on interpretable summaries of the optimization state, including campaign history, estimated hypervolume improvement (EHVI), information gain, and remaining budget and time. By explicitly separating policy selection from candidate generation, the proposed approach enables adaptive control over the exploration–exploitation trade-off while preserving the statistical efficiency and robustness of the underlying optimization framework. 

We evaluate this approach on two discrete alloy design problems formulated as multi-objective black-box optimization tasks over compositional spaces, and assess performance under simulated mid-campaign disruptions such as budget reductions and shortened timelines. By comparing against an exploitation-focused qEHVI baseline and a deterministic fixed mixed allocation strategy, we evaluate whether agent-guided allocation improves campaign responsiveness under evolving constraints. More broadly, this work establishes a generalizable framework for integrating probabilistic optimization with adaptive, language-driven decision-making, and highlights a pathway toward more flexible, scalable, and resilient self-driving laboratory systems operating under realistic, evolving constraints.


